\chapter{Superficies}\label{ch:b2:surfaces}

Tras las curvas, las superficies: objetos de dos parámetros en $\R^3$.
El cálculo diferencial del~\cref{ch:b2:diffcalc} proporciona todo lo
que necesitamos: las derivadas parciales dan los vectores tangentes, el
producto vectorial da la normal y los \omterm{def:b2:linalg:det}{determinantes} dan las
\omterm{def:b2:surfaces:area}{áreas}. Definimos las \omterm{def:b2:surfaces:param}{superficies parametrizadas} regulares, sus
\omterm{def:b2:surfaces:tangent}{planos tangentes} y la \emph{\omterm{def:b2:surfaces:fff}{primera forma fundamental}}, que codifica
todas las medidas de \omterm{def:b2:curves:length}{longitud} y de \omterm{def:b2:surfaces:area}{área} sobre la superficie. Las
superficies también aparecen como conjuntos de nivel $f(x, y, z) = c$;
el gradiente dirige entonces la normal.

\section{Superficies parametrizadas}

\begin{definition}[Superficie parametrizada regular]\label{def:b2:surfaces:param}
Sea $U \subseteq \R^2$ un \omterm{def:b2:metric:topology}{abierto}. Una \emph{superficie
parametrizada}\index{superficie parametrizada} de clase
$\mathcal{C}^k$ ($k \geq 1$) es una aplicación $\sigma \colon U \to
\R^3$, $(u, v) \mapsto \sigma(u, v)$, de clase $\mathcal{C}^k$. Un
punto es \emph{regular} si los vectores de derivadas parciales
\[
\sigma_u(u, v) = \frac{\partial\sigma}{\partial u}(u,v),
\qquad
\sigma_v(u, v) = \frac{\partial\sigma}{\partial v}(u,v)
\]
son linealmente independientes, es decir, si $\sigma_u \wedge
\sigma_v \neq 0$; la superficie es \emph{regular} si lo es todo
punto.
\end{definition}

\begin{example}[Las tres descripciones estándar]\label{ex:b2:surfaces:standard}
\leavevmode
\begin{enumerate}
\item \emph{Grafo:} $\sigma(u, v) = (u,\ v,\ f(u, v))$ para $f \in
\mathcal{C}^1(U)$. Siempre regular: $\sigma_u = (1, 0, f_u)$ y
$\sigma_v = (0, 1, f_v)$ son independientes.
\item \emph{Esfera (coordenadas esféricas):} para la esfera de radio
$R$,
\[
\sigma(\theta, \varphi)
= (R\cos\theta\cos\varphi,\ R\sin\theta\cos\varphi,\
R\sin\varphi),
\qquad
(\theta, \varphi) \in \R \times
\bigl(-\tfrac\pi2, \tfrac\pi2\bigr),
\]
con $\theta$ la \omterm{def:b2:curves:length}{longitud} y $\varphi$ la latitud. Se comprueba que
$\norm{\sigma_\theta \wedge \sigma_\varphi} = R^2\cos\varphi > 0$:
regular fuera de los polos (que esta carta omite).
\item \emph{Conjunto de nivel:} $S = \{(x,y,z) : f(x, y, z) = c\}$
con $f$ de clase $\mathcal{C}^1$ y $\nabla f \neq 0$ sobre $S$.
Cerca de cada punto, una coordenada puede expresarse como función de
las otras dos por el teorema de la función implícita
(\cref{ch:b2:diffcalc}), de modo que $S$ es localmente un grafo.
\end{enumerate}
\end{example}

\begin{example}[Del conjunto de nivel al grafo]
\label{ex:b2:surfaces:levelgraph}
El teorema de la función implícita del punto 3 merece una ejecución
explícita. Tomemos la esfera $x^2 + y^2 + z^2 = R^2$ cerca de su polo
norte $(0, 0, R)$: allí $\frac{\partial f}{\partial z} = 2z = 2R \neq
0$, y despejar $z$ da la carta de grafo
\[
z = \sqrt{R^2 - x^2 - y^2},
\qquad x^2 + y^2 < R^2 ,
\]
regular en todo su dominio (\omterm{def:b2:metric:topology}{abierto}), incluido el polo que la carta
esférica se dejaba. Cerca de un punto del ecuador como $(R, 0, 0)$,
el mismo teorema despeja $x$ en su lugar ($\frac{\partial f}{\partial
x} = 2R \neq 0$). La regla práctica: una superficie de nivel es un
grafo sobre el plano coordenado \emph{ortogonal a la mayor componente
del gradiente}, y cubriendo la esfera con seis cartas de grafo de ese
tipo se comprueba su regularidad en todas partes sin trigonometría
alguna.
\end{example}

\section{Plano tangente y normal}

\begin{definition}[Plano tangente]\label{def:b2:surfaces:tangent}
Sea $\sigma$ regular en $(u_0, v_0)$, y $M_0 = \sigma(u_0, v_0)$. El
\emph{plano tangente}\index{plano tangente} $T_{M_0}S$ es el plano
que pasa por $M_0$ y tiene dirección
$\operatorname{Vect}(\sigma_u, \sigma_v)$ (con las parciales en
$(u_0, v_0)$). La \emph{normal unitaria}\index{normal unitaria} es
\[
n(u_0, v_0)
= \frac{\sigma_u \wedge \sigma_v}
       {\norm{\sigma_u \wedge \sigma_v}} .
\]
\end{definition}

\begin{proposition}[Los vectores tangentes son vectores
velocidad]\label{prop:b2:surfaces:velocity}
La dirección de $T_{M_0}S$ es exactamente el conjunto de vectores
$\gamma'(0)$, donde $\gamma = \sigma \circ c$ recorre las curvas
$\mathcal{C}^1$ trazadas sobre la superficie que pasan por $M_0$ (es
decir, $c \colon (-\varepsilon, \varepsilon) \to U$ es
$\mathcal{C}^1$ con $c(0) = (u_0, v_0)$).
\end{proposition}

\begin{proof}
Si $c(t) = (u(t), v(t))$, la regla de la cadena
(\cref{ch:b2:diffcalc}) da
\[
\gamma'(0) = u'(0)\,\sigma_u + v'(0)\,\sigma_v
\in \operatorname{Vect}(\sigma_u, \sigma_v).
\]
Recíprocamente, el vector $a\sigma_u + b\sigma_v$ se alcanza con la
curva $c(t) = (u_0 + at,\ v_0 + bt)$, que permanece en el \omterm{def:b2:metric:topology}{abierto} $U$
para $\abs t$ pequeño.
\end{proof}

\begin{example}[El plano tangente del helicoide]
\label{ex:b2:surfaces:helicoidtangent}
Para el helicoide $\sigma(u, v) = (v\cos u,\ v\sin u,\ au)$, en el
punto $\sigma(0, 1) = (1, 0, 0)$:
\[
\sigma_u = (0,\ 1,\ a), \qquad \sigma_v = (1,\ 0,\ 0),
\qquad
\sigma_u \wedge \sigma_v = (0,\ a,\ -1),
\]
de modo que el \omterm{def:b2:surfaces:tangent}{plano tangente} es $a\,y = z$. Contiene toda la
generatriz horizontal $\{(t, 0, 0)\}$ (dirección $\sigma_v$): como
en el cono del~\cref{exo:b2:surfaces:1}, en una superficie reglada por
rectas cada generatriz está dentro del \omterm{def:b2:surfaces:tangent}{plano tangente} a lo largo de
ella. La otra dirección tangente $\sigma_u$ es la velocidad de la
hélice $u \mapsto \sigma(u, 1)$: una carta, dos curvas trazadas y
todo el \omterm{def:b2:surfaces:tangent}{plano tangente} queda \omterm{def:b2:structures:generated}{generado}; la~\cref{prop:b2:surfaces:velocity} en acción.
\end{example}

\begin{proposition}[Normal de una superficie de nivel]\label{prop:b2:surfaces:gradient}
Sea $S = \{f = c\}$ con $f$ de clase $\mathcal{C}^1$ y $\nabla
f(M_0) \neq 0$. Entonces el \omterm{def:b2:surfaces:tangent}{plano tangente} de $S$ en $M_0$ es el
plano que pasa por $M_0$ y es ortogonal a $\nabla f(M_0)$:
\[
T_{M_0}S :\quad
\langle \nabla f(M_0),\ M - M_0\rangle = 0 .
\]
\end{proposition}

\begin{proof}
Para toda curva $\gamma$ trazada sobre $S$ que pase por $M_0$,
$f(\gamma(t)) = c$ idénticamente, de modo que la regla de la cadena
da $\langle \nabla f(M_0), \gamma'(0)\rangle = 0$: todos los vectores
velocidad son ortogonales al gradiente, luego la dirección tangente
está contenida en el plano $\nabla f(M_0)^\perp$. Ambos son
subespacios de dimensión $2$ ---la dirección tangente porque $S$ es
localmente un grafo regular (\cref{ex:b2:surfaces:standard}), y el
complemento ortogonal porque $\nabla f(M_0) \neq 0$---, luego son
iguales.
\end{proof}

\begin{example}\label{ex:b2:surfaces:spheretangent}
Para la esfera $x^2 + y^2 + z^2 = R^2$: $\nabla f = 2(x, y, z)$,
luego el \omterm{def:b2:surfaces:tangent}{plano tangente} en $M_0$ es ortogonal al radio $\vect{OM_0}$;
el hecho clásico de que el radio y el \omterm{def:b2:surfaces:tangent}{plano tangente} son
perpendiculares, con ecuación $\langle M_0, M\rangle = R^2$.
\end{example}

\begin{example}[Plano tangente de un grafo]\label{ex:b2:surfaces:graphtangent}
Para $z = f(x, y)$ en $(x_0, y_0)$: aplicando la~\cref{prop:b2:surfaces:gradient} a $F(x,y,z) = f(x,y) - z$,
\[
z = f(x_0, y_0)
+ f_x(x_0, y_0)(x - x_0)
+ f_y(x_0, y_0)(y - y_0) ,
\]
la parte \omterm{def:b2:affine:subspace}{afín} del desarrollo de Taylor de primer orden: el
\omterm{def:b2:surfaces:tangent}{plano tangente} es el grafo de la \omterm{def:b2:diffcalc:differential}{diferencial}, como debe ser.
\end{example}

\begin{example}[El punto más próximo de una
superficie]\label{ex:b2:surfaces:closest}
¿Qué punto del paraboloide $z = x^2 + y^2$ está más cerca de $P =
(0, 0, 1)$? Minimicemos la distancia al cuadrado sobre la superficie:
con $\rho^2 = x^2 + y^2$,
\[
g(\rho^2) = \rho^2 + (\rho^2 - 1)^2,
\qquad
g'(\rho^2) = 1 + 2(\rho^2 - 1) = 0
\iff \rho^2 = \tfrac12 ,
\]
lo que da la circunferencia de puntos de altura $z = \frac12$ y
distancia $\sqrt{\tfrac12 + \tfrac14} = \frac{\sqrt3}2$. La firma
geométrica de la minimalidad: en un punto $M$ así, el vector
$\vect{MP}$ debe ser \emph{normal} a la superficie; de lo contrario,
deslizarse por una curva trazada con velocidad de componente hacia
$P$ disminuiría la distancia. Comprobación: $\nabla(z - x^2 - y^2) =
(-2x, -2y, 1)$ en $M = (x, y, \tfrac12)$, mientras que $\vect{MP} =
(-x, -y, \tfrac12) = \tfrac12(-2x, -2y, 1)$: paralelos, como se
predijo. La condición de primer orden ``pie de la perpendicular'' es
la misma que guiará los extremos sobre conjuntos de nivel en el~\cref{exo:b2:surfaces:12}.
\end{example}

\begin{example}[Planos tangentes de las cuádricas: la regla de
polarización]\label{ex:b2:surfaces:polarization}
Sean $S : xy + yz + zx = 1$ y $M_0 = (1, 1, 0) \in S$. Aquí $\nabla f
= (y + z,\ x + z,\ x + y)$, luego $\nabla f(M_0) = (1, 1, 2)$ y el
\omterm{def:b2:surfaces:tangent}{plano tangente} es
\[
(x - 1) + (y - 1) + 2z = 0,
\qquad\text{es decir,}\qquad x + y + 2z = 2 .
\]
La misma respuesta sale de la \emph{regla de \omterm{def:b2:quadratic:def}{polarización}} que
generaliza el~\cref{ex:b2:surfaces:spheretangent} y el~\cref{exo:b2:surfaces:2}: en la ecuación de la \omterm{pb:b2:surfaces:1}{cuádrica},
sustitúyase $x^2$ por $x_0x$, y cada producto $xy$ por $\frac{x_0y +
y_0x}2$ (y cíclicamente):
\[
\frac{x_0y + y_0x}2 + \frac{y_0z + z_0y}2 + \frac{z_0x +
x_0z}2 = 1
\;\xrightarrow{\ M_0 = (1,1,0)\ }\;
\frac{x + y}2 + \frac z2 + \frac z2 = 1 ,
\]
que es de nuevo $x + y + 2z = 2$. La regla funciona porque el
$\nabla$ de una \omterm{def:b2:quadratic:def}{forma cuadrática} es la forma bilineal asociada
evaluada contra el punto base: la tangencia a una \omterm{pb:b2:surfaces:1}{cuádrica} es
\omterm{def:b2:quadratic:def}{polarización}, una cara más del~\cref{ch:b2:quadratic}.
\end{example}

\section{La primera forma fundamental}

\begin{definition}[Primera forma fundamental]\label{def:b2:surfaces:fff}
Sea $\sigma \colon U \to \R^3$ una superficie $\mathcal{C}^1$
regular. Su \emph{primera forma fundamental}\index{primera forma
fundamental} en $(u,v)$ es la \omterm{def:b2:quadratic:def}{forma cuadrática} definida positiva
sobre $\R^2$
\[
I(h, k) = \norm{h\,\sigma_u + k\,\sigma_v}^2
= E\,h^2 + 2F\,hk + G\,k^2,
\]
donde
\[
E = \norm{\sigma_u}^2,
\qquad
F = \langle \sigma_u, \sigma_v\rangle,
\qquad
G = \norm{\sigma_v}^2 .
\]
\end{definition}

\begin{remark}
$I$ es la restricción del producto escalar euclídeo ambiente al
\omterm{def:b2:surfaces:tangent}{plano tangente}, leída en la base $(\sigma_u, \sigma_v)$: es definida
positiva precisamente porque $\sigma_u, \sigma_v$ son independientes
(\cref{ch:b2:quadratic}). Toda magnitud métrica sobre la superficie
---longitudes de curvas trazadas, ángulos entre ellas, áreas--- se
calcula solo a partir de $E, F, G$. Dos superficies con los mismos
$E, F, G$ en parámetros adecuados son \emph{isométricas} aunque estén
situadas de manera distinta en el espacio: este es el punto de
partida de la geometría intrínseca.
\end{remark}

\begin{example}[Ángulos entre las curvas
coordenadas]\label{ex:b2:surfaces:coordangle}
La \omterm{def:b2:surfaces:fff}{primera forma fundamental} mide también ángulos: las curvas
coordenadas $u \mapsto \sigma(u, v_0)$ y $v \mapsto \sigma(u_0, v)$
se cortan formando el ángulo $\theta$ con
\[
\cos\theta = \frac{\langle\sigma_u,
\sigma_v\rangle}{\norm{\sigma_u}\,\norm{\sigma_v}}
= \frac{F}{\sqrt{EG}} :
\]
el único coeficiente $F$ decide la ortogonalidad de la red de
parámetros. Para la carta de la esfera y para el helicoide, $F = 0$:
los meridianos cortan a los paralelos, y las hélices a las
generatrices horizontales, en ángulo recto; por eso sus integrandos
de \omterm{def:b2:surfaces:area}{área} colapsaron a $\sqrt{EG}$. Para una carta de grafo, $F =
f_xf_y$ solo se anula donde lo hace alguna derivada parcial: la red
coordenada de un grafo inclinado \emph{no} es ortogonal, aunque la
red $(x, y)$ de abajo sí lo sea. Cuando los cálculos sobre una
superficie se ponen pesados, el primer movimiento es buscar una carta
con $F = 0$.
\end{example}

\begin{example}[La carta de la silla de montar]\label{ex:b2:surfaces:saddlechart}
Para la silla de montar $z = xy$ con la carta $\sigma(u, v) = (u, v,
uv)$:
\[
\sigma_u = (1, 0, v), \qquad \sigma_v = (0, 1, u),
\qquad
E = 1 + v^2, \quad F = uv, \quad G = 1 + u^2 ,
\]
y $EG - F^2 = 1 + u^2 + v^2 > 0$: regular en todas partes. Las dos
curvas coordenadas por un punto son \emph{rectas} de $\R^3$ (fíjese
$u$ o fíjese $v$: las generatrices de la silla doblemente reglada) y,
sin embargo, $F \neq 0$ fuera de los ejes: las generatrices por un
punto genérico no son ortogonales. Ambas generatrices están en el
\omterm{def:b2:surfaces:tangent}{plano tangente}, al que generan, de modo que el \omterm{def:b2:surfaces:tangent}{plano tangente}
corta a la superficie a lo largo de dos rectas enteras, el extremo
opuesto de la esfera, cuyos \omterm{def:b2:surfaces:tangent}{planos tangentes} tocan en un solo punto.
El signo del ``contacto de segundo orden'' entre una superficie y sus
\omterm{def:b2:surfaces:tangent}{planos tangentes} es una historia de \omterm{thm:b2:curves:frenet2d}{curvatura}, retomada en el
volumen del tercer año.
\end{example}

\begin{proposition}[Longitud de una curva trazada sobre una
superficie]\label{prop:b2:surfaces:curvelength}
Si $\gamma(t) = \sigma(u(t), v(t))$, $t \in [a, b]$, es
$\mathcal{C}^1$, entonces
\[
L(\gamma) = \int_a^b
\sqrt{E\,u'^2 + 2F\,u'v' + G\,v'^2}\;\dd t ,
\]
con $E, F, G$ evaluados en $(u(t), v(t))$.
\end{proposition}

\begin{proof}
$\gamma' = u'\sigma_u + v'\sigma_v$ por la regla de la cadena, luego
$\norm{\gamma'}^2 = I(u', v') = Eu'^2 + 2Fu'v' + Gv'^2$; intégrese
$\norm{\gamma'}$ (\cref{def:b2:curves:length}).
\end{proof}

\begin{example}[Por qué los aviones vuelan sobre el polo]\label{ex:b2:surfaces:polarroute}
Dos aeropuertos están a la latitud $\varphi_0$ y en longitudes
opuestas: $A = \sigma(0, \varphi_0)$ y $B = \sigma(\pi, \varphi_0)$
sobre la esfera de radio $R$. A lo largo del paralelo ($\varphi
\equiv \varphi_0$), la \omterm{def:b2:curves:length}{longitud} es $\int_0^\pi
R\cos\varphi_0\,\dd\theta = \pi R\cos\varphi_0$. Por la ruta sobre el
polo (subiendo por el meridiano $\theta = 0$ y bajando por el
meridiano $\theta = \pi$), es $2R(\frac\pi2 - \varphi_0)$. A la
latitud $\varphi_0 = \frac\pi3$ (sesenta grados): ruta por el
paralelo $\pi R/2 \approx 1.571\,R$, ruta polar $\pi R/3 \approx
1.047\,R$; un tercio más corta. De hecho, $\pi\cos\varphi_0 \geq \pi
- 2\varphi_0$ sobre $\intcc0{\pi/2}$ (la función $\pi\cos\varphi -
\pi + 2\varphi$ se anula en ambos extremos y su derivada $2 -
\pi\sin\varphi$ cambia de signo una vez, de modo que primero crece y
luego decrece, y por tanto es no negativa): la ruta polar nunca
pierde. La \omterm{def:b2:surfaces:fff}{primera forma fundamental} convirtió una cuestión de
navegación en dos integrales de una línea; el~\cref{exo:b2:surfaces:6} lleva la idea hasta una demostración
genuina de minimalidad para los meridianos.
\end{example}

\begin{lemma}[Identidad de Lagrange]\label{lem:b2:surfaces:lagrange}
Para todos $a, b \in \R^3$: $\norm{a \wedge b}^2 = \norm a^2 \norm
b^2 - \langle a, b\rangle^2$. En particular,
\[
\norm{\sigma_u \wedge \sigma_v} = \sqrt{EG - F^2}.
\]
\end{lemma}

\begin{proof}
Ambos miembros quedan inalterados si sustituimos $b$ por su
componente $b_\perp = b - \frac{\langle a, b\rangle}{\norm a^2}a$
ortogonal a $a$ (para $a \neq 0$; el caso $a = 0$ es trivial): el
izquierdo porque $a \wedge a = 0$, y el derecho desarrollando
$\norm{b_\perp}^2 = \norm b^2 - \frac{\langle a,b\rangle^2}{\norm
a^2}$ y $\langle a, b_\perp\rangle = 0$. Basta, pues, demostrar la
identidad para $a, b$ ortogonales, donde se lee $\norm{a \wedge b} =
\norm a \norm b$: cierto, ya que para vectores ortogonales el
producto vectorial tiene \omterm{def:b2:nvs:norm}{norma} $\norm a\norm
b\,\abs{\sin\frac\pi2}$. La fórmula mostrada es el caso $a =
\sigma_u$, $b = \sigma_v$.
\end{proof}

\begin{remark}
La identidad de Lagrange dice que $EG - F^2 =
\det\begin{psmallmatrix} E & F\\ F & G\end{psmallmatrix}$ es el
\emph{\omterm{def:b2:linalg:det}{determinante} de Gram} de $(\sigma_u, \sigma_v)$: el cuadrado del
\omterm{def:b2:surfaces:area}{área} del paralelogramo que generan. La regularidad, el carácter
definido positivo de la \omterm{def:b2:surfaces:fff}{primera forma fundamental} y la positividad
del \omterm{def:b2:linalg:det}{determinante} de Gram son tres formulaciones de una misma
condición; por eso el integrando de \omterm{def:b2:surfaces:area}{área} de más abajo nunca se anula
sobre una carta regular.
\end{remark}

\begin{definition}[Área]\label{def:b2:surfaces:area}
Sea $\sigma \colon U \to \R^3$ una superficie $\mathcal{C}^1$
regular e inyectiva, y $K \subseteq U$ un dominio \omterm{def:b2:metric:compact}{compacto} sobre el
que tengan sentido las integrales dobles (\cref{ch:b2:multint}). El
\emph{área}\index{área!de una superficie} del trozo $\sigma(K)$ es
\[
\mathcal{A}
= \iint_K \norm{\sigma_u \wedge \sigma_v}\,\dd u\,\dd v
= \iint_K \sqrt{EG - F^2}\;\dd u\,\dd v .
\]
\end{definition}

\begin{remark}[Por qué esta fórmula]
El rectángulo $[u, u + \dd u] \times [v, v + \dd v]$ se aplica, a
primer orden, sobre el paralelogramo \omterm{def:b2:structures:generated}{generado} por $\sigma_u\,\dd u$ y
$\sigma_v\,\dd v$, cuya \omterm{def:b2:surfaces:area}{área} es $\norm{\sigma_u \wedge
\sigma_v}\,\dd u\,\dd v$: la definición integra el factor local de
distorsión de \omterm{def:b2:surfaces:area}{área}, exactamente como la \omterm{def:b2:curves:length}{longitud de arco} integra la
rapidez local. La coherencia con los \omterm{def:b2:curves:reparam}{cambios de parámetro} es el~\cref{exo:b2:surfaces:8}; la coherencia con la fórmula del cambio de
variables para integrales dobles se discute en el~\cref{ch:b2:multint}.
\end{remark}

\begin{example}[Dos grafos distintos, una misma
área]\label{ex:b2:surfaces:samearea}
Sobre el disco unidad, comparemos el cuenco $z = \frac12(x^2 + y^2)$
y la silla de montar $z = xy$. Sus integrandos de \omterm{def:b2:surfaces:area}{área}
(\cref{exo:b2:surfaces:5}) son
\[
\sqrt{1 + x^2 + y^2}
\qquad\text{y}\qquad
\sqrt{1 + y^2 + x^2} :
\]
idénticos. Las dos superficies ---una que se curva igual en todas las
direcciones, la otra con forma de silla--- tienen \omterm{def:b2:surfaces:area}{áreas} exactamente
iguales sobre todo dominio, $\frac{2\pi}3(2\sqrt2 - 1)$ sobre el
disco unidad. El elemento de \omterm{def:b2:surfaces:area}{área} solo ve la \emph{\omterm{def:b2:curves:length}{longitud}} del
gradiente, no la disposición de la flexión; distinguir el cuenco de
la silla exige datos de segundo orden (la estructura de signos
exhibida en la~\cref{fig:b2:surfaces:saddle}), que ninguna medida de
\omterm{def:b2:surfaces:area}{área} detecta. \omterm{def:b2:surfaces:fff}{Primera forma fundamental}: métrica, ciega a la
forma; la segunda forma, que sí ve la forma, pertenece al tercer
año.
\end{example}

\begin{example}[Área de la esfera]\label{ex:b2:surfaces:spherearea}
Para la carta esférica del~\cref{ex:b2:surfaces:standard}:
\[
\sigma_\theta = R(-\sin\theta\cos\varphi,\
\cos\theta\cos\varphi,\ 0),
\qquad
\sigma_\varphi = R(-\cos\theta\sin\varphi,\
-\sin\theta\sin\varphi,\ \cos\varphi),
\]
luego $E = R^2\cos^2\varphi$, $F = 0$, $G = R^2$ y $\sqrt{EG - F^2} =
R^2\cos\varphi$. De ahí,
\[
\mathcal{A}
= \int_{-\pi/2}^{\pi/2}\!\!\int_0^{2\pi}
R^2\cos\varphi\;\dd\theta\,\dd\varphi
= 2\pi R^2\,\bigl[\sin\varphi\bigr]_{-\pi/2}^{\pi/2}
= \boxed{4\pi R^2} .
\]
\end{example}

\begin{example}[El cono, contrastado con la fórmula
escolar]\label{ex:b2:surfaces:conearea}
Para el cono $z = \sqrt{x^2 + y^2}$ sobre la corona $a \leq \rho \leq
b$, la fórmula para grafos del~\cref{exo:b2:surfaces:5} da $1 + f_x^2
+ f_y^2 = 2$ (calcúlense $f_x = x/\rho$, $f_y = y/\rho$), luego
\[
\mathcal A = \sqrt2\,\pi\,(b^2 - a^2) .
\]
Comprobación de coherencia con la fórmula de la generatriz $\pi\rho\ell$
del~\cref{ex:b2:surfaces:revolution}: los conos completos de radios
de base $b$ y $a$ tienen \omterm{def:b2:surfaces:area}{áreas} laterales $\pi b\cdot b\sqrt2$ y $\pi
a\cdot a\sqrt2$, cuya diferencia es exactamente $\sqrt2\pi(b^2 -
a^2)$. Dos cartas, dos fórmulas, un \omterm{def:b2:surfaces:area}{área}: la invariancia demostrada
en el~\cref{exo:b2:surfaces:8}, vista en libertad.
\end{example}

\begin{remark}[Comprobaciones de sensatez para las áreas]
Tres comprobaciones instantáneas cazan la mayoría de los errores en
un cálculo de \omterm{def:b2:surfaces:area}{área}. \emph{Escalado}: dilatar una superficie por
$\lambda$ multiplica $E, F, G$ por $\lambda^2$ y el \omterm{def:b2:surfaces:area}{área} por
$\lambda^2$; una respuesta cuya dependencia de $R$ no sea cuadrática
(como $4\pi R^2$) está mal. \emph{Positividad del elemento}:
$\sqrt{EG - F^2}$ ha de ser estrictamente positivo en el interior de
la carta; un valor nulo señala una degeneración de la carta, que hay
que extirpar como en los polos de la esfera. \emph{Simetría}: un
cálculo sobre un trozo \omterm{def:b2:quadratic:adjoint}{simétrico} debe ser coherente con sumar sus
partes congruentes; más vale que el hemisferio dé $2\pi R^2$.
\end{remark}

\begin{example}[Superficie de revolución]\label{ex:b2:surfaces:revolution}
Gírese la curva $z \mapsto (r(z), 0, z)$, con $r > 0$ de clase
$\mathcal{C}^1$, en torno al eje $z$:
\[
\sigma(\theta, z) = (r(z)\cos\theta,\ r(z)\sin\theta,\ z).
\]
Entonces $E = r(z)^2$, $F = 0$, $G = 1 + r'(z)^2$, luego
\[
\mathcal{A} = \int_a^b\!\!\int_0^{2\pi}
r(z)\sqrt{1 + r'(z)^2}\;\dd\theta\,\dd z
= 2\pi\int_a^b r(z)\sqrt{1 + r'(z)^2}\,\dd z ,
\]
la fórmula clásica (la circunferencia $2\pi r$ por el elemento de
\omterm{def:b2:curves:length}{longitud} de la generatriz). Para el cono $r(z) = kz$, $z \in [0, h]$:
$\mathcal A = 2\pi k\sqrt{1 + k^2}\,\frac{h^2}{2} = \pi \rho \ell$
con $\rho = kh$ el radio de la base y $\ell = h\sqrt{1 + k^2}$ la
generatriz; la fórmula escolar, ahora deducida en vez de admitida.
\end{example}

\begin{example}[El catenoide]\label{ex:b2:surfaces:catenoid}
Gírese la catenaria $r(z) = \cosh z$, $z \in \intcc{-1}{1}$, en torno
a su eje: el \emph{catenoide} resultante tiene, por la fórmula de
revolución y $1 + \sinh^2 z = \cosh^2 z$,
\[
\mathcal A = 2\pi\int_{-1}^1\cosh z\,\sqrt{1 + \sinh^2z}\,
\dd z = 2\pi\int_{-1}^1\cosh^2z\,\dd z
= \pi\bigl[z + \sinh z\cosh z\bigr]_{-1}^{1},
\]
es decir,
\[
\mathcal A = 2\pi + \pi\sinh 2 \approx 17.68 .
\]
El factor de la generatriz $\sqrt{1 + r'^2}$ se fundió con el radio
en un cuadrado perfecto; la misma identidad que hizo elemental la
\omterm{def:b2:curves:length}{longitud de arco} de la catenaria en el capítulo de curvas. No es un
accidente algebraico: entre todas las superficies de revolución que
abarcan las dos circunferencias de borde, el catenoide minimiza el
\omterm{def:b2:surfaces:area}{área} (es la forma de una película de jabón entre dos aros), y esa
propiedad variacional es precisamente lo que singulariza a $\cosh$;
el volumen del tercer año lo demuestra con el cálculo de
variaciones.
\end{example}

\begin{figure}[ht]
\centering
\begin{tikzpicture}[scale=1.05]
  % A saddle surface z = x^2 - y^2 drawn as a wire mesh in oblique projection
  % projection: (x,y,z) -> (x + 0.45*y, z + 0.35*y)
  \begin{scope}
    % u-lines (x varies, y fixed)
    \foreach \yy in {-1,-0.5,0,0.5,1} {
      \draw[ocDarkBlue!70, thick, domain=-1:1, smooth, samples=25]
        plot ({\x + 0.45*\yy}, {\x*\x - \yy*\yy + 0.35*\yy});
    }
    % v-lines (y varies, x fixed)
    \foreach \xx in {-1,-0.5,0,0.5,1} {
      \draw[ocMint!80!black, thick, domain=-1:1, smooth, samples=25]
        plot ({\xx + 0.45*\x}, {\xx*\xx - \x*\x + 0.35*\x});
    }
    % tangent plane at origin: z = 0 -> parallelogram
    \draw[ocRed, dashed]
      (-0.8 - 0.36, -0.28) -- (0.8 - 0.36, -0.28)
      -- (0.8 + 0.36, 0.28) -- (-0.8 + 0.36, 0.28) -- cycle;
    \fill (0,0) circle (0.045) node[below right, font=\small] {$M_0$};
    % normal vector at origin
    \draw[->, thick] (0,0) -- (0,0.9)
      node[right, font=\small] {$n$};
  \end{scope}
\end{tikzpicture}
\caption{La silla de montar $z = x^2 - y^2$ cerca del origen, con sus
curvas coordenadas (las curvas $u$ en azul, las curvas $v$ en verde),
el \omterm{def:b2:surfaces:tangent}{plano tangente} en $M_0 = (0,0,0)$ (a trazos) y la
\omterm{def:b2:surfaces:tangent}{normal unitaria} $n$. La superficie cruza su \omterm{def:b2:surfaces:tangent}{plano tangente}: el
análogo bidimensional de una inflexión.}
\label{fig:b2:surfaces:saddle}
\end{figure}

\begin{remark}[Errores frecuentes]
(i) \emph{Las singularidades de una carta no son singularidades de la
superficie}: la carta esférica degenera en los polos ($\cos\varphi =
0$), pero la esfera es perfectamente regular allí; otra carta
(intercambiando los papeles de los ejes) es regular en los polos.
Antes de declarar singular un punto, pruébese con una segunda
parametrización. (ii) \emph{La regularidad de $\sigma$ concierne a la
parametrización, no a la imagen}: $\sigma(u, v) = (u^3, v, 0)$ no es
regular sobre $u = 0$ aunque su imagen sea un plano. (iii) La fórmula
del \omterm{def:b2:surfaces:area}{área} exige que $\sigma$ sea \emph{inyectiva} sobre $K$: una
carta que cubra un trozo dos veces lo cuenta dos veces (con $\theta$
recorriendo $\intcc0{4\pi}$ se duplica el \omterm{def:b2:surfaces:area}{área} de la esfera). (iv)
La \omterm{def:b2:surfaces:tangent}{normal unitaria} queda definida salvo signo por la superficie,
pero la \emph{elige} la carta (el orden de $u, v$); los enunciados
sobre orientación deben fijar esa elección. (v) Por último, $EG - F^2
> 0$ no es una hipótesis adicional: es exactamente la regularidad,
por la identidad de Lagrange; si se anula en algún punto, el problema
es la carta, y allí no se aplica ninguna fórmula de \omterm{def:b2:surfaces:area}{área} ni de \omterm{def:b2:surfaces:tangent}{plano
tangente}.
\end{remark}

\begin{remark}[Perspectivas dentro de este volumen]
Enlaces hacia delante desde aquí. El elemento de \omterm{def:b2:surfaces:area}{área} $\sqrt{EG -
F^2}\,\dd u\,\dd v$ es un jacobiano bidimensional disfrazado, y el~\cref{ch:b2:multint} hace exacta la analogía con el teorema del
cambio de variables: las integrales de superficie de allí son las
\omterm{def:b2:surfaces:area}{áreas} de este capítulo con un integrando a bordo. La
\omterm{def:b2:surfaces:fff}{primera forma fundamental} es un campo de \omterm{def:b2:quadratic:def}{formas cuadráticas}
positivas, tratadas punto a punto con las herramientas del~\cref{ch:b2:quadratic}, cuyo teorema espectral impulsa también la
clasificación de \omterm{pb:b2:surfaces:1}{cuádricas} del problema de fin de semana de este
capítulo. Y la recta normal guía los problemas de extremos sobre
conjuntos de restricciones (\cref{ex:b2:surfaces:closest}), el germen
geométrico del método de los multiplicadores de Lagrange esbozado con
el~\cref{thm:b2:diffcalc:extrema}.
\end{remark}

\begin{omfigure}
\begin{tikzpicture}[scale=0.8]
  \draw[omDef, thick] (0,0) ellipse (1.5 and 0.95);
  \draw[omThm] (0,0) ellipse (1.5 and 0.35);
  \node[below, font=\small] at (0,-1.25) {elipsoide};
\end{tikzpicture}
\qquad
\begin{tikzpicture}[scale=0.8]
  \draw[omDef, thick] plot[domain=-1.1:1.1, samples=30]
    ({0.6*(exp(\x)+exp(-\x))/2}, {\x});
  \draw[omDef, thick] plot[domain=-1.1:1.1, samples=30]
    ({-0.6*(exp(\x)+exp(-\x))/2}, {\x});
  \draw[omThm] (0,0) ellipse (0.6 and 0.16);
  \draw[omDef] (0,1.1) ellipse (1.0 and 0.22);
  \draw[omDef] (0,-1.1) ellipse (1.0 and 0.22);
  \node[below, font=\small] at (0,-1.55)
    {hiperboloide de una hoja};
\end{tikzpicture}
\qquad
\begin{tikzpicture}[scale=0.8]
  \draw[omDef, thick] plot[domain=-1.05:1.05, samples=30]
    ({\x}, {1.1*\x*\x - 1.1});
  \draw[omThm] (0,0.055) ellipse (1.0 and 0.22);
  \node[below, font=\small] at (0,-1.45)
    {paraboloide elíptico};
\end{tikzpicture}
\qquad
\begin{tikzpicture}[scale=0.8]
  \draw[omDef, thick] plot[domain=-1:1, samples=30]
    ({\x}, {0.55*\x*\x});
  \draw[omThm, thick] plot[domain=-1:1, samples=30]
    ({0.45*\x}, {-0.55*\x*\x + 0.18*\x});
  \node[below, font=\small] at (0,-1.15)
    {paraboloide hiperbólico};
\end{tikzpicture}

{\small Cuatro de las nueve superficies \omterm{pb:b2:surfaces:1}{cuádricas} clasificadas en el
problema de fin de semana, esbozadas mediante sus siluetas y una
curva de nivel (en rojo): el elipsoide acotado, el hiperboloide de
una hoja doblemente reglado con su garganta, el cuenco del
paraboloide elíptico y la silla de montar, cuyas dos secciones
parabólicas se doblan en sentidos opuestos.}
\end{omfigure}

\begin{remark}[Dónde se usa]
El elemento de \omterm{def:b2:surfaces:area}{área} $\norm{\sigma_u\wedge\sigma_v}\,\dd u\,\dd v$ es
la medida de las integrales de superficie del~\cref{ch:b2:multint},
donde se encuentra con la fórmula de Green; la
\omterm{def:b2:surfaces:fff}{primera forma fundamental} es el prototipo de un campo de
\omterm{def:b2:quadratic:def}{formas cuadráticas}, estudiadas punto a punto con las herramientas
del~\cref{ch:b2:quadratic}; y el problema de fin de semana de este
capítulo clasifica todas las superficies \omterm{pb:b2:surfaces:1}{cuádricas} con el teorema
espectral. El volumen del tercer año vuelve a las superficies con las
formas \omterm{def:b2:diffcalc:differential}{diferenciales} y el teorema de la divergencia, y la
\omterm{thm:b2:curves:frenet2d}{curvatura} intrínseca ---lo que $E, F, G$ saben sobre la flexión--- es
la puerta de entrada a la geometría diferencial propiamente dicha.
\end{remark}

\section{Ejercicios}

\begin{exercise}[$\star$]\label{exo:b2:surfaces:1}
Prueba que los \omterm{def:b2:surfaces:tangent}{planos tangentes} del cono $z = \sqrt{x^2 + y^2}$
(menos su vértice) pasan todos por el vértice. \emph{(Parametriza
mediante $\sigma(\theta, r) = (r\cos\theta, r\sin\theta, r)$, $r >
0$.)}
\end{exercise}

\begin{exercise}[$\star$]\label{exo:b2:surfaces:2}
Halla el \omterm{def:b2:surfaces:tangent}{plano tangente} del elipsoide $\frac{x^2}{a^2} +
\frac{y^2}{b^2} + \frac{z^2}{c^2} = 1$ en un punto $(x_0, y_0, z_0)$
de la superficie.
\end{exercise}

\begin{exercise}[$\star$]\label{exo:b2:surfaces:3}
Calcula $E, F, G$ para el helicoide $\sigma(u, v) = (v\cos u,\ v\sin
u,\ au)$, $a > 0$, y el \omterm{def:b2:surfaces:area}{área} del trozo $0 \leq u \leq 2\pi$, $0 \leq
v \leq 1$, en forma de integral (evalúala usando $\int\sqrt{v^2 +
a^2}\,\dd v = \frac12\bigl(v\sqrt{v^2+a^2} + a^2\ln(v + \sqrt{v^2 +
a^2})\bigr) + C$).
\end{exercise}

\begin{exercise}[$\star\star$]\label{exo:b2:surfaces:4}
(Toro) Parametriza el toro obtenido al girar la circunferencia de
centro $(R, 0, 0)$ y radio $r < R$ del plano $xz$ en torno al eje
$z$:
\[
\sigma(\theta, \psi) = \bigl((R + r\cos\psi)\cos\theta,\
(R + r\cos\psi)\sin\theta,\ r\sin\psi\bigr).
\]
Calcula $E, F, G$, comprueba la regularidad y prueba que el \omterm{def:b2:surfaces:area}{área} es
$4\pi^2 R r$ (Pappus: la circunferencia media $2\pi R$ por la
\omterm{def:b2:curves:length}{longitud} $2\pi r$ de la circunferencia).
\end{exercise}

\begin{exercise}[$\star\star$]\label{exo:b2:surfaces:5}
Prueba que el \omterm{def:b2:surfaces:area}{área} del grafo de $f \in \mathcal{C}^1(K)$ es $\iint_K
\sqrt{1 + f_x^2 + f_y^2}\,\dd x\,\dd y$, y calcúlala para el trozo de
paraboloide $z = \frac12(x^2 + y^2)$ sobre el disco $x^2 + y^2 \leq
1$ (coordenadas polares, \cref{ch:b2:multint}).
\end{exercise}

\begin{exercise}[$\star\star$]\label{exo:b2:surfaces:6}
Una curva trazada $\gamma(t) = \sigma(u(t), v(t))$ sobre la esfera de
radio $R$ (carta esférica) tiene $u = \theta(t)$, $v = \varphi(t)$.
Escribe su \omterm{def:b2:curves:length}{longitud} como integral en $\theta, \varphi$ y demuestra
que, entre las curvas que unen dos puntos de un mismo meridiano
$\theta = \theta_0$, el arco de meridiano es la más corta.
\emph{(Acota el integrando inferiormente por $R\abs{\varphi'}$.)}
\end{exercise}

\begin{exercise}[$\star\star\star$]\label{exo:b2:surfaces:7}
(Rectas normales de una esfera) Sea $S$ una superficie de nivel
regular $\{f = c\}$, \omterm{def:b2:metric:connected}{conexa}, cuyas rectas normales pasan todas por
un punto fijo $\Omega$. Prueba que $S$ está contenida en una esfera
centrada en $\Omega$. \emph{(Prueba que $\norm{M - \Omega}^2$ tiene
derivada nula a lo largo de toda curva trazada sobre $S$.)}
\end{exercise}

\begin{exercise}[$\star\star\star$]\label{exo:b2:surfaces:8}
(El \omterm{def:b2:surfaces:area}{área} es geométrica) Sea $\Phi \colon U' \to U$ un difeomorfismo
$\mathcal{C}^1$ entre \omterm{def:b2:metric:topology}{abiertos} de $\R^2$ y sea $\tilde\sigma = \sigma
\circ \Phi$. Prueba que
\[
\norm{\tilde\sigma_{u'} \wedge \tilde\sigma_{v'}}
= \abs{\det J_\Phi}\,
\norm{(\sigma_u \wedge \sigma_v)\circ\Phi} ,
\]
y deduce, usando la fórmula del cambio de variables del~\cref{ch:b2:multint}, que el \omterm{def:b2:surfaces:area}{área} de la~\cref{def:b2:surfaces:area} no depende de la parametrización regular
elegida.
\end{exercise}

\begin{exercise}[$\star$]\label{exo:b2:surfaces:9}
(Teorema de la caja de sombreros de Arquímedes) Sobre la esfera de
radio $R$, la \emph{zona} comprendida entre las latitudes con $z_1
\leq z \leq z_2$ ($-R \leq z_1 < z_2 \leq R$) tiene \omterm{def:b2:surfaces:area}{área} $2\pi
R\,(z_2 - z_1)$: demuéstralo con la carta esférica y concluye que el
\omterm{def:b2:surfaces:area}{área} de una zona depende solo de su altura; cortar una naranja en
rodajas de igual grosor da cantidades iguales de piel.
\end{exercise}

\begin{exercise}[$\star\star$]\label{exo:b2:surfaces:10}
Prueba que toda recta normal de una superficie de revolución
$\sigma(\theta, z) = (r(z)\cos\theta,\ r(z)\sin\theta,\ z)$ ($r > 0$
de clase $\mathcal C^1$) corta al eje de revolución, y localiza el
punto de intersección.
\end{exercise}

\begin{exercise}[$\star\star$]\label{exo:b2:surfaces:11}
(Desenrollar el cilindro) La carta $\sigma(u, v) = (\cos u,\ \sin
u,\ v)$ del cilindro unidad tiene $E = G = 1$, $F = 0$: compruébalo y
explica por qué toda curva trazada $t \mapsto \sigma(u(t), v(t))$
tiene la misma \omterm{def:b2:curves:length}{longitud} que la curva plana $t \mapsto (u(t),
v(t))$. Deduce que la hélice que va de $(1, 0, 0)$ a $(1, 0, 2\pi c)$
dando una vuelta tiene \omterm{def:b2:curves:length}{longitud} $2\pi\sqrt{1 + c^2}$, y que ninguna
curva trazada con los mismos extremos y una vuelta completa es más
corta.
\end{exercise}

\begin{exercise}[$\star\star\star$]\label{exo:b2:surfaces:12}
Sea $S = \{f = c\}$ una superficie de nivel regular \omterm{def:b2:metric:compact}{compacta} y $M_0
\in S$ un punto a distancia máxima del origen. Prueba que $\nabla
f(M_0)$ es colineal con $\vect{OM_0}$: la normal en el punto más
lejano es radial. Aplícalo al elipsoide $\frac{x^2}{a^2} +
\frac{y^2}{b^2} + \frac{z^2}{c^2} = 1$ ($a > b > c > 0$): halla todos
los puntos donde la normal es radial e identifica los más lejanos.
\end{exercise}

\section{Problema: la clasificación de las cuádricas de $\R^3$}

\begin{problem}[{Problema de fin de semana --- todas las superficies
cuádricas, ordenadas por el teorema espectral}]\label{pb:b2:surfaces:1}
Una \emph{cuádrica}\index{cuádrica} es el conjunto de ceros en $\R^3$
de un polinomio de grado dos
\[
q(X) = X^{\mathsf T}\!AX + 2\,\langle b, X\rangle + c,
\qquad A \in \mathcal S_3(\R),\ A \neq 0,\ b \in \R^3,\
c \in \R .
\]
Las superficies de las figuras de este capítulo ---esferas,
elipsoides, sillas de montar, conos, cilindros--- son todas
\omterm{pb:b2:surfaces:1}{cuádricas}. Este problema las clasifica \emph{todas}: el teorema
espectral (\cref{thm:b2:quadratic:spectral}) endereza la parte
cuadrática, las traslaciones \omterm{def:b2:affine:subspace}{afines} (\cref{ch:b2:affine}) absorben la
parte lineal, y lo que queda es una lista corta y completa de formas
normales.

\medskip
\noindent\textbf{Parte I --- La máquina de reducción.}
\begin{enumerate}
  \item Sea $X = PY + t$ con $P \in O(3)$ y $t \in \R^3$ (un cambio
        rígido de coordenadas). Prueba que $q(PY + t) = Y^{\mathsf
        T}\!A'Y + 2\langle b', Y\rangle + c'$ con
        \[
        A' = P^{\mathsf T}\!AP, \qquad
        b' = P^{\mathsf T}(At + b), \qquad
        c' = q(t) .
        \]
        Deduce que el \omterm{def:b2:reduction:eigen}{espectro} de $A$ (y por tanto su rango y su
        signatura) es un invariante rígido de la ecuación, y explica
        por qué la ecuación de una \omterm{pb:b2:surfaces:1}{cuádrica} dada solo está
        determinada salvo un factor escalar no nulo.
  \item Usando el teorema espectral, prueba que tras una rotación la
        ecuación pasa a ser $\sum_i \lambda_i y_i^2 +
        2\sum_i\beta_iy_i + c = 0$ con $\lambda_1, \lambda_2,
        \lambda_3$ los \omterm{def:b2:reduction:eigen}{valores propios} de $A$.
  \item Para cada $i$ con $\lambda_i \neq 0$, absorbe $\beta_iy_i$
        mediante una traslación ($y_i \mapsto y_i -
        \beta_i/\lambda_i$). Escribe la ecuación reducida cuando
        $\operatorname{rank} A = r$: $\sum_{i\leq r}\lambda_i z_i^2 +
        2\sum_{i > r}\beta_iz_i + c'' = 0$.
  \item Un \emph{centro} de la \omterm{pb:b2:surfaces:1}{cuádrica} de ecuación $q = 0$ es un
        punto $\Omega$ con $q(2\Omega - X) = q(X)$ para todo $X$: la
        simetría central en $\Omega$ conserva la ecuación y, por
        tanto, la superficie. Prueba que $q(2\Omega - X) - q(X) =
        -4\langle A\Omega + b, X\rangle + 4\langle A\Omega + b,
        \Omega\rangle$, y deduce que los centros son exactamente las
        soluciones de $A\Omega = -b$; existen si y solo si $b \in
        \operatorname{im}A$, y el centro es único si y solo si $A$ es
        invertible.
\end{enumerate}

\medskip
\noindent\textbf{Parte II --- \omterm{pb:b2:surfaces:1}{Cuádricas} con centro
($\operatorname{rank}A = 3$).} Aquí la ecuación reducida es
$\lambda_1z_1^2 + \lambda_2z_2^2 + \lambda_3z_3^2 = \delta$.
\begin{enumerate}[resume]
  \item Multiplicando por $-1$ si hace falta, supóngase que al menos
        dos $\lambda_i > 0$. Enumera las posibilidades: signatura
        $(3, 0)$ con $\delta > 0$, $= 0$, $< 0$; y signatura $(2, 1)$
        con $\delta > 0$, $= 0$, $< 0$; nombra los seis conjuntos
        resultantes (elipsoide, punto, conjunto vacío, hiperboloide
        de una hoja, cono, hiperboloide de dos hojas) y pon cada uno
        en su forma normal euclídea ($\frac{x^2}{a^2} +
        \frac{y^2}{b^2} + \frac{z^2}{c^2} = 1$, etc.).
  \item Clasifica $x^2 + y^2 + z^2 + 4xy + 4yz + 4zx = 1$: prueba que
        $A = 2J - I$ con $J$ la matriz de unos, calcula el \omterm{def:b2:reduction:eigen}{espectro}
        $\{5, -1, -1\}$ e identifica un hiperboloide de dos hojas de
        revolución en torno al eje $\R(1,1,1)$.
  \item (Generatrices) Para el hiperboloide de una hoja
        $\frac{x^2}{a^2} + \frac{y^2}{b^2} - \frac{z^2}{c^2} = 1$,
        factoriza
        \[
        \Bigl(\frac xa - \frac zc\Bigr)
        \Bigl(\frac xa + \frac zc\Bigr)
        = \Bigl(1 - \frac yb\Bigr)\Bigl(1 + \frac yb\Bigr)
        \]
        y produce dos familias uniparamétricas de rectas contenidas
        en la superficie.
  \item Prueba que por cada punto del hiperboloide de una hoja pasa
        exactamente una recta de cada familia: la superficie es
        \emph{doblemente reglada}.
  \item El \emph{cono asintótico} del hiperboloide de una hoja es $C
        : \frac{x^2}{a^2} + \frac{y^2}{b^2} - \frac{z^2}{c^2} = 0$.
        Con $\rho^2 = \frac{x^2}{a^2} + \frac{y^2}{b^2}$, prueba que
        todo punto del hiperboloide está a distancia a lo sumo
        $c\bigl(\rho - \sqrt{\rho^2 - 1}\bigr) = \frac{c}{\rho +
        \sqrt{\rho^2 - 1}}$ de $C$, de modo que la superficie se
        abraza a su cono en el infinito. ¿Cuáles son las secciones
        del hiperboloide por los planos $x = \pm a$?
\end{enumerate}

\medskip
\noindent\textbf{Parte III --- Rango $2$ y rango $1$: paraboloides,
cilindros, planos.}
\begin{enumerate}[resume]
  \item Supongamos $\operatorname{rank}A = 2$, digamos $\lambda_1,
        \lambda_2 \neq 0 = \lambda_3$. Partiendo de la pregunta 3,
        sepárense dos casos según sea $\beta_3 \neq 0$ (sin centro,
        por la pregunta 4) o $\beta_3 = 0$ (una recta de centros), y
        redúzcase a
        \[
        \lambda_1z_1^2 + \lambda_2z_2^2 + 2\beta_3z_3 = 0
        \qquad\text{o}\qquad
        \lambda_1z_1^2 + \lambda_2z_2^2 + c'' = 0 :
        \]
        paraboloides elípticos o hiperbólicos en el primer caso, y
        cilindros sobre cónicas con centro (o pares de planos que se
        cortan, una recta, el conjunto vacío) en el segundo.
  \item Prueba que la silla de montar $z = xy$ es un paraboloide
        hiperbólico: gírese $\pi/4$ en el plano $xy$ para llegar a $z
        = \tfrac12(u^2 - v^2)$, la superficie de la~\cref{fig:b2:surfaces:saddle} salvo escala.
  \item Prueba que la silla de montar $z = xy$ lleva las dos familias
        de rectas $\{x = x_0,\ z = x_0y\}$ y $\{y = y_0,\ z =
        xy_0\}$, con exactamente una recta de cada familia por cada
        punto: la segunda \omterm{pb:b2:surfaces:1}{cuádrica} doblemente reglada.
  \item Clasifica $x^2 + y^2 - 2x + 4y + 3 = 0$ en $\R^3$ (complétense
        los cuadrados; identifica un cilindro circular recto y da su
        eje y su radio).
  \item Sea ahora $\operatorname{rank}A = 1$, digamos $\lambda_1 \neq
        0 = \lambda_2 = \lambda_3$. Girando dentro del plano núcleo y
        trasladando, redúzcase a
        \[
        \lambda_1z_1^2 + 2\beta z_2 = 0 \quad (\beta \neq 0)
        \qquad\text{o}\qquad
        \lambda_1z_1^2 + c'' = 0 :
        \]
        un cilindro parabólico, o bien un par de planos paralelos, un
        plano doble o el conjunto vacío. Clasifica $(x + y)^2 = z$
        por completo (forma normal, eje de invariancia por
        traslación).
\end{enumerate}

\medskip
\noindent\textbf{Parte IV --- El teorema de clasificación.}
\begin{enumerate}[resume]
  \item Ensambla las partes I--III en un teorema: \emph{toda
        \omterm{pb:b2:surfaces:1}{cuádrica} de $\R^3$ se lleva mediante un movimiento rígido a
        exactamente una forma normal}. Enumera los diecisiete tipos
        \omterm{def:b2:affine:subspace}{afines} (cuenta las variantes vacías y los conjuntos
        degenerados) y destaca las nueve \emph{superficies}
        \omterm{pb:b2:surfaces:1}{cuádricas}: elipsoide, hiperboloides de una y de dos hojas,
        cono, paraboloides elíptico e hiperbólico, y cilindros
        elíptico, hiperbólico y parabólico.
  \item Escribe el \emph{algoritmo} de clasificación: dados $(A, b,
        c)$, ¿qué magnitudes calculas, en qué orden y qué rama decide
        qué tipo? Justifica que cada paso es efectivo
        (\omterm{def:b2:reduction:eigen}{valores propios} de una matriz \omterm{def:b2:quadratic:adjoint}{simétrica} $3\times3$, rango,
        resolubilidad de $A\Omega = -b$).
  \item Ejecuta el algoritmo sobre $x^2 + y^2 + z^2 - 2xy - 2yz -
        2zx = 1$: prueba que $A = 2I - J$ tiene \omterm{def:b2:reduction:eigen}{espectro} $\{2, 2,
        -1\}$ y concluye: un hiperboloide de una hoja de revolución
        en torno a $\R(1,1,1)$.
  \item Ejecútalo sobre $x^2 + y^2 - z^2 - 2x + 4y + 2z + 4 = 0$:
        halla el centro e identifica la \omterm{pb:b2:surfaces:1}{cuádrica}.
  \item Ejecútalo sobre $x^2 + 4xy + y^2 = 2z$: diagonaliza el bloque
        $xy$ ($u = \frac{x+y}{\sqrt2}$, $v = \frac{x-y}{\sqrt2}$) e
        identifica la \omterm{pb:b2:surfaces:1}{cuádrica}.
  \item Euclídeo frente a \omterm{def:b2:affine:subspace}{afín}. Prueba que dos \omterm{pb:b2:surfaces:1}{cuádricas} con centro
        en forma normal son \emph{rígidamente} equivalentes si y solo
        si tienen las mismas listas de coeficientes (salvo
        permutación y un escalar positivo común sobre la ecuación),
        mientras que \emph{afínmente} solo sobreviven los datos de la
        signatura: todo elipsoide es imagen \omterm{def:b2:affine:subspace}{afín} de la esfera
        redonda. ¿Qué teorema garantiza que la signatura no puede
        cambiar por el camino (\cref{thm:b2:quadratic:sylvester})?
\end{enumerate}

\medskip
\noindent\textbf{Parte V --- Dividendos.}
\begin{enumerate}[resume]
  \item Prueba que toda sección de una \omterm{pb:b2:surfaces:1}{cuádrica} por un plano \omterm{def:b2:affine:subspace}{afín} es
        una cónica (posiblemente degenerada) de ese plano. Identifica
        las secciones de la silla de montar $z = xy$ por los planos
        $z = c$ ($c \neq 0$ y $c = 0$).
  \item ¿Qué superficies \omterm{pb:b2:surfaces:1}{cuádricas} contienen rectas? Prueba que el
        elipsoide, el hiperboloide de dos hojas y el paraboloide
        elíptico no contienen ninguna \emph{(restringe $q$ a una
        recta y usa la desigualdad de Cauchy--Schwarz para el caso de
        dos hojas)}; que el cono y los cilindros están reglados por
        una familia; y concluye que las superficies \omterm{pb:b2:surfaces:1}{cuádricas}
        doblemente regladas son exactamente el hiperboloide de una
        hoja y el paraboloide hiperbólico.
  \item Cuando solo se busca el tipo \emph{\omterm{def:b2:affine:subspace}{afín}}, la \omterm{thm:b2:quadratic:gauss}{reducción de
        Gauss} (\cref{thm:b2:quadratic:gauss}) sale más barata que
        diagonalizar. Rehaz la pregunta 17 con el algoritmo de Gauss
        y comprueba la signatura $(2, 1)$; ¿qué información euclídea
        pierde Gauss?
  \item Todos los \omterm{def:b2:reduction:eigen}{valores propios} de una matriz \omterm{def:b2:quadratic:adjoint}{simétrica} son reales;
        prueba que, en consecuencia, los \emph{signos} de los
        \omterm{def:b2:reduction:eigen}{valores propios} de $A$ pueden leerse en el
        \omterm{def:b2:reduction:charpoly}{polinomio característico} mediante la regla de los signos de
        Descartes, y verifícalo sobre la pregunta 6:
        $\chi_A(\lambda) = \lambda^3 - 3\lambda^2 - 9\lambda - 5$
        tiene exactamente un cambio de signo, luego signatura $(1,
        2)$.
  \item Síntesis. Resume el algoritmo en unas pocas líneas; enuncia
        el papel exacto que desempeñan (i) el teorema espectral, (ii)
        la ecuación del centro $A\Omega = -b$, (iii) el teorema de
        inercia de Sylvester y (iv) la \omterm{thm:b2:quadratic:gauss}{reducción de Gauss}. ¿Qué da la
        misma máquina en $\R^2$, y qué cambia en $\R^n$?
\end{enumerate}
\end{problem}
